\documentclass[10pt,a4paper]{amsart}
\usepackage[margin=19mm]{geometry}
\usepackage{amsmath,amssymb,mathtools}
\usepackage[scr=rsfs,cal=euler]{mathalfa}
\usepackage{xfrac}
\usepackage[unicode,hidelinks,bookmarks=false]{hyperref}
\newcommand{\ie}{i.e.\ }
\newcommand{\cf}{cf.\ }
\newcommand{\resp}{respectively\ }
\newcommand{\Subsection}{\subsection}
\providecommand{\nobreakdash}{\nobreak\hbox{-}}
\setlength{\emergencystretch}{2em}
\multlinegap0pt
\usepackage{algorithm}\usepackage{algpseudocode}
\newtheorem{theorem}{Theorem}
\newcommand{\M}[1]{\mathbf{#1}}
\newcommand{\ptl}{\partial}
\newcommand{\til}[1]{\widetilde{#1}}
\title{From Walking to Tunneling: An Investigation of Generalized Pilot-Wave Dynamics}
\author{Akilan Sankaran, Diego Israel Chavez}
\date{Source: arXiv:2512.20110v1 (CC BY 4.0)}
\begin{document}
\maketitle

\noindent\emph{Source and licence.} From Walking to Tunneling: An Investigation of Generalized Pilot-Wave Dynamics. Version arXiv:2512.20110v1 (CC BY 4.0), distributed by arXiv under the Creative Commons Attribution 4.0 International licence, http://creativecommons.org/licenses/by/4.0/, as recorded in the arXiv licence metadata for this exact version and shown on the abstract page https://arxiv.org/abs/2512.20110v1. Full licence notice: \texttt{licenses/arxiv-2512.20110-CC-BY-4.0.txt}.\par\smallskip

\noindent\textit{Editorial scope.} Counted unit: the article's theorem nondim\_procedure (source0.tex lines 267--283). The article's complete original proof of it is delivered with it (source0.tex lines 285--332, one \texttt{proof} environment of 48 lines). No mathematics is written, repaired or re-derived: the statement and the proof are the article's own lines, in the article's own notation, and no retained line is substituted or re-flowed.  Retained uncounted as the source-local context the proof needs: lines 248--258.  Nothing was omitted from any retained span, so the package carries no omission record.  No retained line contains a citation, so the wrapper carries no bibliography and no bibliography entry.  No retained line contains a figure environment, an image-inclusion command or drawing code, so no image is delivered and the package declares no asset dependency.  Editorial formatting only, in addition: an AMS \texttt{amsart} preamble that restates the article's own declarations and macros used by the retained lines, namely the environments \texttt{theorem} on separate counters, together with the standard packages those lines need.  The retained environments print in this document as Theorem 1 in order of appearance, whereas the article prints the same items with the numbering of its own full document.  The complete native source of the article as distributed in the arXiv e-print for this exact version is bundled unchanged as \texttt{sources/191537/source0.tex}.
    \begin{align}
\centering
\label{thm_3_first_first_eq}
    0 &= \nabla \phi, \\
\label{bottom_thm_3.1}
    0 &= \phi_z - \phi_x H_x - \phi_y H_y,\\
\label{phi_t_thm_3.1}
    \phi_t &= -g(t) \eta + \frac{\sigma}{\rho} \Delta_{H} \eta  + 2 \nu \Delta_{H} \phi - \frac{1}{\rho} P_D(\M{x} - \M{x}_p(t), t) ,\\
\label{eta_t_thm_3.1}
    \eta_t &= \phi_z + 2 \nu \Delta_{H} \eta,
\end{align}

\begin{theorem}
\label{nondim_procedure}
The evolution of the wave-field and a walking droplet on a fluid bath with variable topography can be characterized through the following differential equations:
\begin{align}
    \label{laplacian_1}
    0 &= \mu^2 \Delta_{{H}}{{\phi}} + \phi_{zz}, \\
\label{bottom_1}
    0 &= \phi_z - \mu^2({{\phi}_{{x}} {H}_{{x}} + {\phi}_{{y}} {H}_{{y}}}),\\
\label{weird_1}
    \phi_t &= - G \cdot  (1 + \gamma \cos(\omega t)) {\eta} + \frac{2}{\mathrm{Re}}  \Delta_{{H}}{{\phi}} + \mathrm{Bo} \, \Delta_{{H}}{{\eta}} - \frac{GM}{\rho} P_D({\M{{x}}} - {\M{{x}}_p(t)}, t), \\
\label{kinematical_condition_1}
    {\eta}_{{t}} &=  \frac{1} {\mu} {\phi}_{{z}} +\frac{2}{\mathrm{Re}} \Delta_H{{\eta}}, \\
\label{position_update}
    - F(t) \nabla \eta (\M{x}_p) &=  m \M{x}_p'' + c F(t) \M{x}_p'.
\end{align}
with\footnote{For completeness, we note that $T_F$ designates the Faraday period, and $\lambda_F$ designates the Faraday wavelength.} $\mu = \frac{h}{\lambda_F}, G = \frac{\text{g} T_F^2}{\lambda_F}$, $\mathrm{Bo} = \frac{\sigma T_F^2}{\rho \lambda_F^3}$, $M = \frac{m}{\rho \lambda_F^3}$.
\end{theorem}

\begin{proof}[Proof.] We introduce the following nondimensionalization schema:
\begin{align*}
    \Bigg \{
    \begin{array}{ll}
      \til{x} = \ell x, \, \, \, \til{y} = \ell y, \, \, \, \til{z} = hy, \, \, \, \\
    \til{t} = t_F t, \, \, \, \til{\eta} = \eta_c \eta, \, \, \, \til{\phi} = \phi_c \phi,\til{H} = H_c H. 
    \end{array}
\end{align*}
 
where $H_c, \eta_c, \ell$ are constants and $h, t_F$ give the mean depth and Faraday period. We now apply this schema to equations (\ref{thm_3_first_first_eq}–\ref{eta_t_thm_3.1}). For equation (\ref{thm_3_first_first_eq}), we have:
\begin{align*}
    \frac{\ptl}{\ptl x} \left(\phi_c \frac{\ptl \til{\phi}}{\ptl \til{x}} \frac{1}{\ell} \right) + \frac{\ptl}{\ptl y} \left(\phi_c \frac{\ptl \til{\phi}}{\ptl \til{y}} \frac{1}{\ell} \right) + \frac{\ptl}{\ptl z} \left(\phi_c \frac{\ptl \til{\phi}}{\ptl \til{z}} \frac{1}{h} \right) = 0.
\end{align*}

Simplification and rearrangement in accordance with the horizontal Laplacian gives the following elliptic partial differential equation, with $\mu := h/\ell$ being a nondimensional parameter:
\begin{align}
\label{nondim-laplace}
    \mu^2 \Delta_{\til{H}}{\til{\phi}} + \frac{\ptl^2 \til{\phi}}{\ptl \til{z}^2} = 0.
\end{align}

Similarly, we may express the nondimensionalized form of equation (\ref{bottom_thm_3.1}) as follows:
\begin{align*}
    \phi_c \frac{\ptl \til{\phi}}{\ptl \til{z}} \frac{1}{z_c} - \phi_c  \frac{\ptl \til{\phi}}{\ptl \til{x}} \frac{1}{\ell} \cdot H_c \frac{\ptl \til{H}}{\ptl \til{x}} \frac{1}{\ell}  - \phi_c  \frac{\ptl \til{\phi}}{\ptl \til{y}} \frac{1}{\ell} \cdot H_c \frac{\ptl \til{H}}{\ptl \til{y}} \frac{1}{\ell} = 0.
\end{align*}
Applying our substitution for the nondimensional parameter $\mu$ gives the following:
\begin{align}
\label{nondim-bottom}
    \frac{\ptl \til{\phi}}{\ptl \til{z}} - \mu^2{\til{\phi}_{\til{x}} \til{H}_{\til{x}} + \til{\phi}_{\til{y}} \til{H}_{\til{y}}} = 0.
\end{align}

Next, we nondimensionalize the stress condition (\ref{phi_t_thm_3.1}) as follows:
\begin{align*}
    \phi_c \frac{\ptl \til{\phi}}{\ptl \til{t}} \frac{1}{t_F} &= -g(t) \eta_{c} \til{\eta} + \frac{\sigma}{\rho} \cdot \frac{\eta_c^2}{x_c^2} \left( \frac{\ptl^2 \til{\eta}}{\ptl \til{x}^2} + \frac{\ptl^2 \til{\eta}}{\ptl \til{y}^2} \right) \\
    & \, \, \, + 2 \nu \frac{\phi_c^2}{x_c^2} \Delta_{\til{H}}{\til{\phi}} - \frac{1}{\rho} P_D(x_c ({\M{\til{x}}} - {\M{\til{x}}_p(t)}), t_c \til{t}).
\end{align*}
After including the Reynolds number Re := $\frac{2 \nu t_F}{\lambda_F^2}$ for our fluid, we may write:
\begin{align}
\label{eqn33}
    \frac{\phi_c}{t_F} \frac{\ptl \til{\phi}}{\ptl \til{t}} &= -g(t) \eta_{c} \til{\eta} + \mathrm{Re}  \Delta_{\til{H}}{\til{\nu}} + 2 \nu \frac{\phi_c^2}{x_c^2} \Delta_{\til{H}}{\til{\phi}} - \frac{1}{\rho} P_D(x_c ({\M{\til{x}}} - {\M{\til{x}}_p(t)}), T_F \til{t}).
\end{align}

Finally, the analogue of the kinematic condition (\ref{eta_t_thm_3.1}) in the nondimensional case is:
\begin{align}
\label{eqn44}
    \til{\eta}_{\til{t}} = \frac{t_F \phi_c}{\eta_c h} \til{\phi}_{\til{z}} + \mathrm{Re} \nabla_{\til{H}}{\til{\eta}}.
\end{align}
Equations (\ref{nondim-laplace}-\ref{eqn44}) specify the nondimensionalized version of our system.
\end{proof}

\end{document}
